\chapter*{Chapter 1 glossary}\addcontentsline{toc}{chapter}{Chapter 1 glossary}
\begin{description}
\item[Stored structure] The persistent description available for future activity. See Section~\ref{sec:stored}.
\item[Expressed computation] The ordered activity selected or constructed for a context. See Section~\ref{sec:stored}.
\item[Adapted structure] A persistent description after a specified update process. See Section~\ref{sec:stored}.
\item[Gene expression] Production of functional RNA or protein from genomic information. See Section~\ref{sec:biology}.
\item[Regulation] Biological control of activity; in the proposed software interface, a context-dependent control decision. See Section~\ref{sec:biology}.
\item[Plan] An ordered list of operations with defined composition semantics. See Section~\ref{sec:example}.
\item[Development] Biological construction processes; by explicit analogy, a software description-to-program map. See Section~\ref{sec:development}.
\item[State migration] A declared way to move or reset state when structure changes. See Section~\ref{sec:timescales}.
\item[Novelty burden] The need to identify a contribution beyond established alternatives. See Section~\ref{sec:alternatives}.
\item[DOGMA] The non-Transformer DNA-native model and corresponding engine research line. See Section~\ref{sec:taxonomy}.
\item[Hermon DNA] The Transformer-based DNA model and corresponding engine research line. See Section~\ref{sec:taxonomy}.
\item[Evolutor] The broader research, compiler/runtime and orchestration framework above both families. See Section~\ref{sec:taxonomy}.
\item[KV cache] Stored key/value representations for retained context positions. See Section~\ref{sec:memory}.
\item[Parity] Agreement with a specified reference behavior under declared tolerances. See Section~\ref{sec:roadmap}.
\end{description}
